\section{Arquivo ATP}

O arquivo \codefile{surto\_bobina.atp} descreve um caso de transitório eletromagnético no formato ATP. O ATP é uma variante do EMTP voltada a simulações no domínio do tempo em sistemas elétricos, incluindo elementos RLC, fontes, chaves, linhas e equipamentos de potência \cite{atpdrawhelp}. A origem histórica do EMTP está associada ao trabalho de Dommel e ao desenvolvimento na Bonneville Power Administration \cite{dommel1969,emtp_history}. Programas EMTP são amplamente usados para surtos de manobra, descargas atmosféricas, coordenação de isolamento e outros eventos rápidos \cite{long1990}.

\paragraph{Nota de revisão.} O deck originalmente versionado neste projeto \emph{não executava}: a primeira tentativa real (abril/2026) terminou com erro fatal \codefile{KILL = 6} porque os cartões de ramo estavam desalinhados em relação ao formato de colunas fixas do ATP. A auditoria técnica do projeto (\codefile{AUDITORIA\_PROJETO.md}) identificou ainda dois erros latentes: a unidade de indutância (com \codefile{XOPT = 0} o ATP lê mH, e não henries como o comentário original afirmava — o valor antigo seria interpretado 1000 vezes menor) e o tipo da fonte (cartão tipo 14, que no ATP é fonte cossenoidal, e não dupla exponencial). O deck foi corrigido e executado com sucesso em 12/06/2026; as evidências da falha original e da execução corrigida estão preservadas em \codefile{docs/evidencias\_atp\_falha/}.

O cabeçalho do arquivo documenta o caso: bobina distribuída, impulso 1.2/50~\si{\micro\second}, escada Pi com \Nsec{} seções, \(L_{\mathrm{total}}=\SI{10}{mH}\), \(R_{\mathrm{total}}=\SI{5}{\ohm}\), \(C_{\mathrm{total}}=\SI{1}{nF}\), fonte dupla exponencial de \SI{1}{kV} de pico e terminação aberta. Com \codefile{XOPT = 0} as indutâncias dos cartões de ramo são interpretadas em \si{mH} e com \codefile{COPT = 0} as capacitâncias em \si{\micro F}.

A seção de dados de simulação define \(\Delta t=\SI{1e-8}{s}\) e \(T_{\max}=\SI{2e-4}{s}\). Esse tempo máximo é maior que o \codefile{t\_total} do caso Python principal (\SI{50}{\micro\second}), o que permite observar mais ciclos de reflexão. No cartão inteiro de dados diversos, \codefile{IOUT = 1000} limita a impressão tabular no \codefile{.lis}, \codefile{IPLOT = 1} grava o arquivo de plotagem a cada passo — a grade temporal do \codefile{.pl4} (\SI{1e-8}{s}) coincide exatamente com a grade de reporte do Python — e \codefile{ICAT = 1} mantém o \codefile{.pl4} em disco após a execução (sem ele, o GNU ATP descarta o arquivo ao final). A malha R--L possui 20 ramos série de \SI{0.25}{\ohm} e \SI{0.5}{mH} cada. Os capacitores shunt são de \SI{5e-5}{\micro F} (\SI{50}{pF}) nos nós internos e \SI{2.5e-5}{\micro F} (\SI{25}{pF}) nos extremos.

A fonte é o cartão tipo 15 do ATP (função de surto dupla exponencial), com a forma
\begin{equation}
V(t) = A_1\left(e^{A t} - e^{B t}\right),
\qquad A < 0,\; B < 0,
\end{equation}
em que os expoentes são fornecidos no cartão como valores \emph{negativos}: \(A_1=\SI{1035.1}{V}\), \(A=\SI{-13863}{s^{-1}}\) e \(B=\SI{-2.5e6}{s^{-1}}\). O fator \(A_1\) incorpora a normalização para que o pico fique próximo de \SI{1}{kV}. A interpretação do cartão foi verificada de duas formas: pelo eco de interpretação do \codefile{.lis} (\codefile{Source.~1.04E+03 -1.39E+04 -2.50E+06}) e pela coincidência, dentro do arredondamento de ponto flutuante, entre a tensão \(V(N0)\) registrada no \codefile{.pl4} e a expressão analítica acima. A seção de saída solicita as tensões nos nós \(N0\), \(N5\), \(N10\), \(N15\) e \(N20\), correspondendo à entrada, 25\%, 50\%, 75\% e saída aberta da bobina.

Com a execução real do deck, o ATP deixa de ser apenas uma referência conceitual: ele fornece uma solução independente do mesmo circuito (método trapezoidal de integração, na tradição EMTP \cite{dommel1969}), contra a qual o modelo Python é comparado quantitativamente na Seção~\ref{sec:correspondencia}.
